\documentclass[11pt,a4paper]{article}

\usepackage[a4paper,margin=2.3cm]{geometry}
\usepackage[T1]{fontenc}
\usepackage[utf8]{inputenc}
\usepackage{amsmath}
\usepackage{amssymb}
\usepackage{booktabs}
\usepackage{array}
\usepackage{parskip}
\usepackage{hyperref}
\usepackage{pifont}

\newcommand{\code}[1]{\texttt{#1}}

\title{Uncertainty breakdown of the combined top-quark mass:\\ three impact definitions}
\author{}
\date{}

\begin{document}
\maketitle
\vspace{-2.5em}

\begin{quote}
\textit{Math note for the two-stage combination (pyconvino cross-section
combination $\to$ mtpole-ttj-pyconvino $\chi^2$ mass fit). Covers (A) the
classic freeze/conditional impact, (B) the marginal ``global impact'' of
arXiv:2307.04007, and (C) the exact additive variance-share generalisation used
for correlated priors.}
\end{quote}

All numbers quoted are \textbf{uncertainties only} (the fit is blinded); they come
from the 3-experiment \code{Combination\_ATLAS813CMS13\_corrV2} fit (ATLAS 8~TeV +
ATLAS 13~TeV + CMS 13~TeV, NNLO / CT18NNLO, \code{polyOrder 1}).

\section*{0. Setup and notation}

\textbf{Stage 1 --- combination.} \code{pyconvino} combines the input differential
measurements into a vector of combined cross sections $x\in\mathbb{R}^{n}$ (here
$n=19$ bins), profiling $N$ experimental nuisance parameters
$\theta\in\mathbb{R}^{N}$ (here $N=614$). Each nuisance carries a Gaussian prior

\[
\theta \sim \mathcal{N}\!\big(0,\ C\big),\qquad C^{-1}\equiv \code{inv\_C}\ (\text{prior precision}).
\]

For the usual ``unit, independent'' priors $C=\mathbb{1}$; the \code{corrV2} setup
deliberately makes $C$ (hence \code{inv\_C}) \textbf{non-diagonal}, correlating
common systematics \emph{across experiments}. The combination returns the joint
post-fit covariance

\[
\text{cov\_full}=
\begin{pmatrix}\Sigma_{xx} & \Sigma_{x\theta}\\[2pt] \Sigma_{\theta x} & \Sigma_{\theta\theta}\end{pmatrix},
\qquad
\Sigma_{xx}\equiv C_{\rm tot}=\code{combined\_covariance},\quad
\Sigma_{x\theta}=\code{x\_sys\_cov}.
\]

\textbf{Stage 2 --- mass fit.} In the Gaussian/linear regime the mass estimator is
a linear functional of the combined cross sections,

\[
\boxed{\ \hat m - m_0 = w^\top\,(x-x_0)\ }\qquad
w=\text{Fisher weights (\code{\_fisher\_weights})}.
\]

(The code applies a normalisation Jacobian $J$ so that $w_{\rm abs}=J^\top w$ acts
in absolute-cross-section space; that is an implementation detail and is
suppressed below.)

\textbf{The one vector that drives everything.} Define the covariance of the mass
with each nuisance,

\[
\boxed{\ g_i \equiv \mathrm{Cov}(\hat m,\theta_i) = \big(w^\top \Sigma_{x\theta}\big)_i\ }
\qquad (\text{code: } \code{dm\_i} = w_{\rm abs}\!\cdot\!\code{x\_sys\_cov}[:,i]).
\]

Also write $\sigma_{{\rm post},i}=\sqrt{\Sigma_{\theta\theta,ii}}$ (post-fit width
of $\theta_i$) and $\sigma_{{\rm pre},i}=\sqrt{C_{ii}}$ (prior width). The total
experimental mass variance is $\sigma_m^2 = w^\top C_{\rm tot}\,w$.

Everything below is a different way of attributing $\sigma_m^2$ (or its
systematic part) to the individual $\theta_i$.

\section*{A. Classic impact --- freeze / conditional method}

\textbf{Recipe.} Fix one nuisance to its best fit $\pm$ its \textbf{post-fit}
uncertainty, $\theta_i=\hat\theta_i\pm\sigma_{{\rm post},i}$, re-minimise, read the
shift in $\hat m$.

\textbf{Closed form.} For a Gaussian joint posterior of $(\hat m,\theta)$ the
conditional mean gives

\[
\boxed{\ I^{\rm freeze}_i=\frac{\mathrm{Cov}(\hat m,\theta_i)}{\sigma_{{\rm post},i}}
=\frac{g_i}{\sigma_{{\rm post},i}}=\sigma_m\,\rho(\hat m,\theta_i)\ }
\]

i.e.\ the familiar ``impact $=\sigma_m\times$ correlation''.

\textbf{Stat/syst split (Tier-1).} The systematic is defined by \emph{freezing
all} nuisances:

\[
\sigma_{\rm stat}^{\rm freeze}=\big(\text{mass unc.\ with all }\theta\text{ fixed}\big),\qquad
\big(\sigma_{\rm syst}^{\rm freeze}\big)^2=\sigma_m^2-\big(\sigma_{\rm stat}^{\rm freeze}\big)^2
= w^\top V_{\rm syst}\,w,
\]

with $V_{\rm syst}=C_{\rm tot}-V_{\rm stat}$ (code:
\texttt{total\_\allowbreak syst\_\allowbreak covariance}).
Equivalently $f=\big(w^\top V_{\rm syst} w\big)/\big(w^\top C_{\rm tot} w\big)$ and
$\sigma_{\rm syst}=\sigma_m\sqrt{f}$, $\sigma_{\rm stat}=\sigma_m\sqrt{1-f}$.

\textbf{Key property / caveat.} The freeze \emph{stat/syst split closes}
($\sigma_{\rm stat}\oplus\sigma_{\rm syst}=\sigma_m$), \textbf{but the
per-nuisance $I^{\rm freeze}_i$ do NOT add up} to $\sigma_{\rm syst}^{\rm freeze}$
--- the classic ``impacts don't sum'' fact, because the $\theta_i$ are correlated
in the post-fit and each $I^{\rm freeze}_i$ double-counts shared constraint.

\textit{Live numbers:} $\sigma_{\rm stat}^{\rm freeze}=424.7$~MeV,
$\sigma_{\rm syst}^{\rm freeze}=733.1$~MeV ($\oplus=847.2$~MeV experimental total).

\section*{B. Global impact --- marginal (arXiv:2307.04007)}

\textbf{Recipe.} Shift the \emph{external/auxiliary observable} $y_i$ that
constrains $\theta_i$ by its \textbf{prefit} width $\sigma_{{\rm pre},i}$, let all
$\theta$ \textbf{float}, refit, read $\Delta\hat m$.

\textbf{Relation to the freeze impact.} The two differ by the \emph{post-fit
nuisance-parameter covariance} factor $\sigma_{{\rm post},i}/\sigma_{{\rm pre},i}$:

\[
\boxed{\ I^{\rm glob}_i = I^{\rm freeze}_i\times\frac{\sigma_{{\rm post},i}}{\sigma_{{\rm pre},i}}
=\frac{g_i}{\sigma_{{\rm pre},i}}\ }\ \xrightarrow{\ \text{unit priors }\sigma_{\rm pre}=1\ }\
I^{\rm glob}_i = g_i = \mathrm{Cov}(\hat m,\theta_i).
\]

This $I^{\rm glob}_i=g_i$ is exactly the code's marginal \code{dm\_i} and the
marginal forest bars.

\textbf{Additive property (the reason the method exists).} \emph{When the
auxiliary observables are independent} ($C$, hence \code{inv\_C}, diagonal),

\[
\big(\sigma_{\rm syst}^{\rm glob}\big)^2=\sum_i \big(I^{\rm glob}_i\big)^2=\sum_i g_i^2 .
\]

\textbf{Stat/syst interpretation (paper slide 8).} The global split books the
\emph{in-situ statistical constraint of the nuisances} as \textbf{statistical}, so
$\sigma_{\rm syst}^{\rm glob}=$ fluctuations of the external measurements only,
and $\sigma_{\rm stat}^{\rm glob}=\sqrt{\sigma_m^2-(\sigma_{\rm syst}^{\rm glob})^2}$
absorbs that in-situ constraint. Same total, different stat/syst line than the
freeze method --- neither is ``more correct''.

\textbf{Caveat under correlated priors.} With \code{inv\_C} non-diagonal (our
\code{corrV2} case) the plain quadrature $\sum_i g_i^2$ is \textbf{not}
$\sigma_{\rm syst}^2$ --- see \S~C.

\textit{Live numbers:} $\sqrt{\sum_i g_i^2}=954.7$~MeV, which even
\textbf{exceeds} the experimental total (852.1~MeV) --- a clear symptom that
quadrature-summing the marginals is wrong here.

\section*{C. Global impact --- exact additive variance share (correlated priors)}

\textbf{The correct systematic variance.} Propagate fluctuations of the auxiliary
data $y$ (covariance = prior covariance $C=\code{inv\_C}^{-1}$) through the
profiled fit. Using $\partial \hat x/\partial y=\Sigma_{x\theta}\,\code{inv\_C}$,

\[
\frac{\partial\hat m}{\partial y}=w^\top\Sigma_{x\theta}\,\code{inv\_C}=g^\top\code{inv\_C},
\]

\[
\boxed{\ \sigma_{\rm syst}^2=\frac{\partial\hat m}{\partial y}\,C\,\frac{\partial\hat m}{\partial y}^{\!\top}
=\big(g^\top\code{inv\_C}\big)\,\code{inv\_C}^{-1}\,\big(\code{inv\_C}\,g\big)
= g^\top\,\code{inv\_C}\,g\ }
\]

This is positive ($\code{inv\_C}\succ0$) and provably $\le\sigma_m^2$; it reduces
to $\sum_i g_i^2$ only when $\code{inv\_C}=\mathbb{1}$.

\textbf{Exact additive attribution.} Define $a\equiv\code{inv\_C}\,g$. Then

\[
\sigma_{\rm syst}^2=g^\top a=\sum_i \underbrace{g_i\,a_i}_{\textstyle s_i},\qquad
s_g=\sum_{i\in g} g_i\,a_i,\qquad
\boxed{\ \sum_i s_i=\sum_g s_g=\sigma_{\rm syst}^2\ (\text{exactly})\ }
\]

Plot the signed root $\operatorname{sign}(s)\sqrt{|s|}$ so the (signed) squares
still add to $\sigma_{\rm syst}^2$.

\textbf{Why this is the paper's method, generalised.} With independent priors
$\code{inv\_C}=\mathbb{1}\Rightarrow a_i=g_i\Rightarrow s_i=g_i^2=\big(I^{\rm glob}_i\big)^2$
--- the two agree term by term. Restoring the correlation metric \code{inv\_C} is
the \emph{only} change; it turns the oblique (correlated) axes into the right
inner product.

\textbf{Geometric picture --- why the marginals (\S~B) don't sum but these do.}
$\sigma_{\rm syst}$ is the \emph{length} of $g$ in the \code{inv\_C} metric.
Cholesky $\code{inv\_C}=LL^\top$ and $h=L^\top g$ give an \textbf{orthogonal}
(whitened) basis where Pythagoras holds exactly,
$\sum_k h_k^2=g^\top\code{inv\_C}\,g$ (verified to $10^{-13}$). The named
nuisances are the \emph{oblique} axes; summing the squares of their projections
$g_i^2$ ignores the cross terms $\sum_{i\ne j} g_i(\code{inv\_C})_{ij}g_j$. The
variance share $s_i=g_i a_i$ folds each row's full cross term back in, restoring
closure.

\textbf{Why a share can be negative.}
$s_i<0\iff\operatorname{sign}(g_i)\ne\operatorname{sign}(a_i)$: the source's own
mass pull $g_i$ opposes the correlation-weighted pull
$a_i=(\code{inv\_C}\,g)_i$, i.e.\ it is \emph{anti-correlated through the prior}
with the dominant systematics, so including it slightly \textbf{reduces} the
total variance --- the exact analogue of a negative covariance term in an
ordinary correlated error budget.

\textit{Live numbers:} $\sqrt{g^\top\code{inv\_C}\,g}=830.3$~MeV (linear)
$\le 852.1$ total; per-group shares close to $\sigma_{\rm syst}=825.5$~MeV
(HESSE-anchored) with \textbf{1 of 49} groups negative
(\code{ATLAS\_13TeV\_modelling\_singletop}, $-1.9$~MeV: $g=-1.7$, $a=+2.2$);
per-nuisance shares close to the same, 23 of 614 negative.

\section*{D. Side-by-side}

\begingroup
\renewcommand{\arraystretch}{1.5}
\small
\begin{tabular}{@{}p{2.4cm}p{3.7cm}p{3.7cm}p{5.3cm}@{}}
\toprule
 & \textbf{A. Freeze / classic} & \textbf{B. Global --- marginal} & \textbf{C. Global --- additive share} \\
\midrule
Per-source formula & $I^{\rm freeze}_i=g_i/\sigma_{{\rm post},i}$ & $I^{\rm glob}_i=g_i/\sigma_{{\rm pre},i}$ & $s_i=g_i\,(\code{inv\_C}\,g)_i$ \\
Reads as & ``shift if $\theta_i$ fixed to $\pm\sigma_{\rm post}$'' & ``shift if aux obs moves $\pm\sigma_{\rm pre}$'' & ``contribution of $\theta_i$ to $\sigma_{\rm syst}^2$'' \\
$\sum$ over sources $=\sigma_{\rm syst}^2$? & \ding{55}~(never) & \ding{51}~only if priors independent & \ding{51}~\textbf{always (exact)} \\
Sign & $\ge 0$ & $\ge 0$ & signed (can be $<0$) \\
Books in-situ NP constraint as & systematic & statistical & (attribution of the global $\sigma_{\rm syst}$) \\
Total exp.\ syst & $733.1$~MeV & (marginals overshoot: $954.7$) & $825.5$~MeV \\
Code output & \texttt{doExp\allowbreak Stat\allowbreak Syst\allowbreak Breakdown} (Tier-1) & \texttt{global\_\allowbreak impacts\_\allowbreak\{groups,\allowbreak nuisances\}} & \texttt{global\_\allowbreak impacts\_\allowbreak\{groups,\allowbreak nuisances\}\_\allowbreak additive} \\
\bottomrule
\end{tabular}
\endgroup

\bigskip
\textbf{Relations:}
$I^{\rm glob}_i = I^{\rm freeze}_i\cdot\sigma_{{\rm post},i}/\sigma_{{\rm pre},i}$;
and $s_i \xrightarrow{\ \code{inv\_C}\to\mathbb{1}\ } (I^{\rm glob}_i)^2$. Both
stat/syst splits close to the same experimental total ($847.2$~MeV here); they
only draw the stat/syst line in different places.

\section*{E. Practical guidance}

\begin{itemize}
  \item \textbf{Headline stat vs syst:} either split is defensible; quote which
    convention (freeze vs global) is used. Freeze $=$ traditional; global $=$
    ``external-only systematic''.
  \item \textbf{``Which sources dominate?''} --- read the \textbf{marginal} bars
    (\S~B): honest ``size of each source on its own''. Do \textbf{not}
    quadrature-sum them under correlated priors.
  \item \textbf{``Budget that adds up to $\sigma_{\rm syst}$''} --- use the
    \textbf{additive shares} (\S~C). Expect a few small negative entries
    (anti-correlated sources); that is correct.
  \item Everything here is the \emph{experimental} systematic (in the
    combined-cross-section covariance). Theory (PDF, scale, interpolation) is
    added separately:
    $\sigma_{\rm total}^2=\sigma_{\rm exp}^2+\sigma_{\rm PDF}^2+\sigma_{\rm interp}^2+\sigma_{\rm scale}^2$.
\end{itemize}

\bigskip
\noindent\textit{Reference:} A.~Freitas et al., \emph{``Uncertainty breakdown via
global impacts''}, arXiv:2307.04007 (CMS-week slides by A.~Gilbert, slides
4--9). \S~A $\equiv$ conditional method (slide~3); \S~B $\equiv$ global impact /
marginal additive form (slides~5--6, independent-aux case); \S~C $\equiv$ its
exact quadratic-form generalisation for correlated auxiliary observables.

\end{document}
